\origpage{32}{ROBERTSON AND SEYMOUR}
If $H$ is a drawing in $\Sigma$, a \textit{chain} is a sequence $v_0, r_1, \dots, r_k, v_k$ such that $r_1, \dots, r_k$ are regions of $H$ and for $1 \le i \le k$, $\{v_{i-1}, v_i\} \subseteq \bar{r}_i$. (We do not insist that $v_0, \dots, v_k$ are vertices.) The chain has \textit{length} $k$, and $v_0, v_k$ are its \textit{ends}. If $X, Y \subseteq \Sigma$ are non-empty, we define the \textit{regional distance} between $X$ and $Y$ to be the minimum length of a chain with one end in $X$ and the other in $Y$. From (5.2) we deduce

\stmt
(5.3) \textit{Let $H$ be a drawing in $\Sigma$, and let $\Delta \subseteq \Sigma$ be a closed disc. Let $Z \subseteq \Delta - \text{bd}(\Delta)$ be non-empty and arc-wise connected, such that the regional distance between $Z$ and $\Sigma - \Delta$ is $\ge k \ge 2$. Then there is a circuit $C$ of $H$ with $U(C) \subseteq \Delta$, bounding an open disc $\Lambda$ including $Z$, such that the regional distance between $\bar{\Lambda}$ and $\Sigma - \Delta$ is $\ge k - 1$.}

\textit{Proof.} If $r$ is a region of $H$ and $\bar{r} \cap Z \ne \emptyset$ then $r \subseteq \Delta - \text{bd}(\Delta)$; for otherwise $\bar{r} \cap (\Sigma - \Delta) \ne \emptyset$ and the regional distance between $Z$ and $\Sigma - \Delta$ is $\le 1$, a contradiction. Hence (5.2) may be applied. Let $C$ be as in (5.2), with $U(C)$ bounding an open disc $\Lambda \subseteq \Delta$. We claim that the regional distance between $\bar{\Lambda}$ and $\Sigma - \Delta$ is $\ge k - 1$. For let $v_0, r_1, v_1, \dots, r_{k'}, v_{k'}$ be a chain with $v_0 \in \bar{\Lambda}$ and $v_{k'} \in \Sigma - \Delta$, and with $k'$ minimum. Now $k' \ge 1$ since $\bar{\Lambda} \subseteq \Delta$, and $v_1 \notin \bar{\Lambda}$ since $k'$ is minimum. Hence $r_1 \not\subseteq \Lambda$, and so $r_1 \cap \Lambda = \emptyset$, and $v_0 \in V(C)$. Let $e$ be an edge of $C$ incident with $v_0$. From the choice of $C$ there is a region $r$ of $H$ incident with $e$ satisfying $\bar{r} \cap Z \ne \emptyset$. Choose $v \in \bar{r} \cap Z$; then from the chain
\begin{equation*}
v, r, v_0, r_1, v_1, \dots, r_{k'}, v_{k'}
\end{equation*}
and the hypothesis that the regional distance between $Z$ and $\Sigma - \Delta$ is at least $k$, we deduce that $k' \ge k - 1$, as required. $\square$

By repeated application of (5.3), we deduce

\stmt
(5.4) \textit{Let $H, \Sigma, \Delta, Z, k$ be as in (5.3). Then there are mutually disjoint circuits $C_1, \dots, C_{k-1}$ of $H$ with $U(C_1), \dots, U(C_{k-1}) \subseteq \Delta$, all bounding open discs including $Z$.}

A similar proof \del{(we omit the details)} yields

\stmt
(5.5) \textit{Let $H$ be a drawing in $\Sigma$, and let $\Delta \subseteq \Sigma$ be a closed disc. Let $Z \subseteq \Delta - \text{bd}(\Delta)$ be non-empty and arc-wise connected, such that the regional distance between $Z$ and $\text{bd}(\Delta)$ is $\ge k \ge 1$. Then there are $k - 1$ mutually disjoint circuits $C_1, \dots, C_{k-1}$ of $H$ with $U(C_1), \dots, U(C_{k-1}) \subseteq \Delta - \text{bd}(\Delta)$, all bounding open discs including $Z$.}

\fix{\textit{Proof of (5.5).} We first show the following analogue of (5.3).

\hspace*{12pt}(1) \textit{If $k \ge 2$, there is a circuit $C$ of $H$ with $U(C) \subseteq \Delta - \text{bd}(\Delta)$, bounding an open disc $\Lambda$ including $Z$, such that $\bar{\Lambda} \subseteq \Delta - \text{bd}(\Delta)$ and the regional distance between $\bar{\Lambda}$ and $\text{bd}(\Delta)$ is $\ge k - 1$.}

For let $r$ be a region of $H$ with $\bar{r} \cap Z \ne \emptyset$. Then $\bar{r} \cap \text{bd}(\Delta) = \emptyset$, for otherwise the regional distance between $Z$ and $\text{bd}(\Delta)$ is $\le 1$. In particular $r$ is disjoint from $\text{bd}(\Delta)$. Moreover $r$ meets $\Delta - \text{bd}(\Delta)$, since $\bar{r}$ meets $Z \subseteq \Delta - \text{bd}(\Delta)$ and $\Delta - \text{bd}(\Delta)$ is open. Since $r$ is connected, it follows that $r \subseteq \Delta$, and so $\bar{r} \subseteq \Delta - \text{bd}(\Delta)$. In particular (5.2) may be applied. Let $C$ be as in (5.2), bounding an open disc $\Lambda \subseteq \Delta$ including $Z$. Every edge $e$ of $C$ is incident with a region $r$ of $H$ with $\bar{r} \cap Z \ne \emptyset$, and so $\bar{e} \subseteq \bar{r} \subseteq \Delta - \text{bd}(\Delta)$. Since every vertex of $C$ is an end of an edge of $C$, it follows that $U(C) \subseteq \Delta - \text{bd}(\Delta)$. Since $\Lambda$ is open and $\Lambda \subseteq \Delta$, we have $\Lambda \cap \text{bd}(\Delta) = \emptyset$, and so $\bar{\Lambda} = \Lambda \cup U(C) \subseteq \Delta - \text{bd}(\Delta)$. Now let $v_0, r_1, v_1, \dots, r_{k'}, v_{k'}$ be a chain with $v_0 \in \bar{\Lambda}$ and $v_{k'} \in \text{bd}(\Delta)$, with $k'$ minimum. Then $k' \ge 1$, since $\bar{\Lambda} \cap \text{bd}(\Delta) = \emptyset$. Also $v_1 \notin \bar{\Lambda}$, by the minimality of $k'$ if $k' \ge 2$, and since $v_1 \in \text{bd}(\Delta)$ if $k' = 1$. Hence $r_1 \not\subseteq \Lambda$, for otherwise $v_1 \in \bar{r}_1 \subseteq \bar{\Lambda}$. Since $r_1$ is connected and disjoint from $U(C)$, it follows that $r_1 \cap \Lambda = \emptyset$, and so $\bar{r}_1 \cap \Lambda = \emptyset$, since $\Lambda$ is open. Thus $v_0 \notin \Lambda$, since $v_0 \in \bar{r}_1$, and therefore $v_0 \in U(C)$. Let $e$ be an edge of $C$ with $v_0 \in \bar{e}$, and let $r$ be a region of $H$ incident with $e$ such that $\bar{r} \cap Z \ne \emptyset$. Choose $v \in \bar{r} \cap Z$. Then $v_0 \in \bar{e} \subseteq \bar{r}$, and so
\begin{equation*}
v, r, v_0, r_1, v_1, \dots, r_{k'}, v_{k'}
\end{equation*}
is a chain of length $k' + 1$ between $Z$ and $\text{bd}(\Delta)$. Hence $k' + 1 \ge k$. This proves (1).

We prove (5.5) by induction on $k$. If $k = 1$ there is nothing to prove, and so we assume that $k \ge 2$. Let $C$ and $\Lambda$ be as in (1). Then $\bar{\Lambda}$ is a closed disc, and so it is non-empty and arc-wise connected. Moreover $\bar{\Lambda} \subseteq \Delta - \text{bd}(\Delta)$, and the regional distance between $\bar{\Lambda}$ and $\text{bd}(\Delta)$ is $\ge k - 1 \ge 1$. From our inductive hypothesis, applied with $\bar{\Lambda}$ in place of $Z$, there are $k - 2$ mutually disjoint circuits $C_2, \dots, C_{k-1}$ of $H$ with $U(C_2), \dots, U(C_{k-1}) \subseteq \Delta - \text{bd}(\Delta)$, bounding open discs $\Lambda_2, \dots, \Lambda_{k-1}$ that include $\bar{\Lambda}$, and hence include $Z$. For $2 \le i \le k - 1$, $U(C) \subseteq \bar{\Lambda} \subseteq \Lambda_i$ and $U(C_i) \cap \Lambda_i = \emptyset$, and so $U(C) \cap U(C_i) = \emptyset$. Setting $C_1 = C$, the circuits $C_1, \dots, C_{k-1}$ satisfy (5.5). $\square$}\fixnote{fix:5.5-proof-added}{The paper omits the proof of (5.5). The argument above is the one used in the formalization.}
